\documentclass[11pt]{article}
\usepackage{amsmath,amssymb,amsthm}
\usepackage[margin=1.1in]{geometry}

\newtheorem{theorem}{Theorem}
\newtheorem{lemma}{Lemma}
\newtheorem{corollary}{Corollary}
\newtheorem{definition}{Definition}
\newtheorem{example}{Example}
\newtheorem{remark}{Remark}

\newcommand{\X}{\mathcal{X}}
\newcommand{\Y}{\mathcal{Y}}
\newcommand{\D}{\mathcal{D}}
\newcommand{\Hyp}{\mathcal{H}}
\newcommand{\err}{L_{\D}}
\newcommand{\emperr}{L_{S}}
\newcommand{\VC}{\mathrm{VCdim}}

\title{Learnability}
\author{Yuval Lavie\\[2pt]
\small following Shalev-Shwartz \& Ben-David,
\emph{Understanding Machine Learning} (2014), Part I}
\date{\today}

\begin{document}
\maketitle

\begin{abstract}
When is a learning problem solvable from finitely many examples, and what
does the answer depend on? This note develops the answer end to end:
the formal model and ERM, why unrestricted learning overfits, PAC
learnability with a complete proof for finite classes, uniform convergence
as the engine behind agnostic learning, the No-Free-Lunch theorem (there
is no universal learner), and the VC dimension, which characterizes
learnability exactly. Companion experiments:
\texttt{learnability\_demo.py} exhibits the memorizer overfitting, the
impossibility of shattering two points with thresholds, and the two
learning rates $1/m$ and $1/\sqrt{m}$ of the Fundamental Theorem.
\end{abstract}

\section{The formal model}

\begin{definition}[The learning problem]
A \emph{domain} $\X$ (instances), a \emph{label set} $\Y = \{0,1\}$, and
an unknown distribution $\D$ over $\X \times \Y$. The learner receives a
\emph{training sample} $S = \bigl((x_1,y_1),\dots,(x_m,y_m)\bigr)$ drawn
i.i.d.\ from $\D$, and must output a \emph{hypothesis}
$h:\X \to \Y$. Its quality is the \emph{true risk} (generalization error)
and its observable proxy, the \emph{empirical risk}:
\begin{equation}
  \err(h) \;=\; \Pr_{(x,y)\sim\D}\bigl[h(x)\ne y\bigr],
  \qquad
  \emperr(h) \;=\; \frac1m\,\bigl|\{\,i : h(x_i)\ne y_i\,\}\bigr| .
  \label{eq:risks}
\end{equation}
\end{definition}

The learner never sees $\D$; everything it knows is $S$. In the
\emph{realizable} setting one additionally assumes a true labeling
function $f$ with $y = f(x)$ and some hypothesis achieving zero risk; the
general (\emph{agnostic}) setting assumes nothing.

\begin{definition}[Empirical Risk Minimization]
Given a \emph{hypothesis class} $\Hyp \subseteq \Y^{\X}$ chosen before
seeing data, the ERM rule outputs
\begin{equation}
  h_S \;\in\; \arg\min_{h\in\Hyp} \emperr(h) .
  \label{eq:erm}
\end{equation}
\end{definition}

\section{Why restrict $\Hyp$ at all: overfitting}

Let ERM range over \emph{all} functions. Consider the memorizer
\begin{equation*}
  h_S(x) =
  \begin{cases}
    y_i & \text{if } x = x_i \text{ for some } i,\\
    0 & \text{otherwise.}
  \end{cases}
\end{equation*}
Its empirical risk is $0$ by construction. But if $\X$ is continuous and
$\D$ has a density, a fresh $x$ almost surely differs from every training
point, so $h_S$ predicts the constant $0$ on all new data --- with, e.g.,
$\err(h_S) = \tfrac12$ when labels are balanced. Perfect memory, zero
learning (first panel of the demo). The conclusion is structural, not a
technicality: \emph{some prior restriction --- an inductive bias --- is
necessary}. Sections 5--6 make this precise (No-Free-Lunch) and quantify
exactly how much restriction is needed (VC dimension).

\section{PAC learnability, and the finite realizable case}

\begin{definition}[PAC learnability, realizable]
$\Hyp$ is \emph{PAC learnable} if there exist a function
$m_{\Hyp}(\varepsilon,\delta)$ and a learner such that for every
$\varepsilon,\delta\in(0,1)$ and every realizable $\D$: given
$m \ge m_{\Hyp}(\varepsilon,\delta)$ i.i.d.\ examples, the output $h_S$
satisfies
\begin{equation}
  \Pr\bigl[\err(h_S) \le \varepsilon\bigr] \;\ge\; 1-\delta .
  \label{eq:pac}
\end{equation}
\end{definition}
\emph{Probably} ($1-\delta$, over the random draw of $S$)
\emph{Approximately} ($\varepsilon$) \emph{Correct}. Both slacks are
unavoidable: a finite sample can be unlucky ($\delta$) and can never pin
$\D$ down exactly ($\varepsilon$).

\begin{theorem}[Finite classes are PAC learnable]
\label{thm:finite}
Every finite $\Hyp$ is PAC learnable by ERM with sample complexity
\begin{equation}
  m_{\Hyp}(\varepsilon,\delta)
  \;\le\;
  \Bigl\lceil \frac{\log\bigl(|\Hyp|/\delta\bigr)}{\varepsilon}
  \Bigr\rceil .
  \label{eq:finitesample}
\end{equation}
\end{theorem}

\begin{proof}
Call $h$ \emph{bad} if $\err(h) > \varepsilon$, and let
$\Hyp_B \subseteq \Hyp$ be the bad hypotheses. A fixed bad $h$ labels one
random example correctly with probability at most $1-\varepsilon$, so it
\emph{survives} the sample (i.e.\ $\emperr(h)=0$) with probability
\begin{equation*}
  \Pr\bigl[\emperr(h) = 0\bigr]
  \le (1-\varepsilon)^m \le e^{-\varepsilon m} .
\end{equation*}
By the union bound, the probability that \emph{some} bad hypothesis
survives is at most $|\Hyp|\, e^{-\varepsilon m}$, which is $\le \delta$
once $m \ge \log(|\Hyp|/\delta)/\varepsilon$. On the complementary event
no bad hypothesis is consistent with $S$; and by realizability the perfect
$h^\ast$ \emph{is} consistent, so ERM \eqref{eq:erm} returns a consistent
hypothesis --- necessarily a good one: $\err(h_S)\le\varepsilon$.
\end{proof}

Read the bound: examples are priced logarithmically in the class size and
linearly in $1/\varepsilon$ --- doubling the class costs one more
$\varepsilon^{-1}\log 2$ worth of data, not double the data.

\section{Agnostic learning via uniform convergence}

Realizability is a strong fiction; drop it. The target now is
\emph{competing with the best in the class}.

\begin{definition}[Agnostic PAC]
$\Hyp$ is \emph{agnostically PAC learnable} if with
$m \ge m_{\Hyp}(\varepsilon,\delta)$ samples, for \emph{every} $\D$,
\begin{equation}
  \Pr\Bigl[\err(h_S) \;\le\; \min_{h\in\Hyp}\err(h) + \varepsilon\Bigr]
  \;\ge\; 1-\delta .
  \label{eq:agnostic}
\end{equation}
\end{definition}

The engine is a property of the \emph{sample}, not the algorithm:

\begin{definition}[$\varepsilon$-representative sample; uniform convergence]
$S$ is \emph{$\varepsilon$-representative} for $\Hyp$ if
$\;|\emperr(h) - \err(h)| \le \varepsilon$ for \emph{all} $h \in \Hyp$
simultaneously. $\Hyp$ has the \emph{uniform convergence property} if
some $m^{\mathrm{UC}}_{\Hyp}(\varepsilon,\delta)$ samples make $S$
$\varepsilon$-representative with probability $\ge 1-\delta$, for every
$\D$.
\end{definition}

\begin{lemma}[UC $\Rightarrow$ agnostic PAC via ERM]
\label{lem:ucerm}
If $S$ is $\tfrac{\varepsilon}{2}$-representative, then the ERM output
satisfies $\err(h_S) \le \min_{h\in\Hyp} \err(h) + \varepsilon$.
\end{lemma}

\begin{proof}
For any $h\in\Hyp$:
$\err(h_S) \le \emperr(h_S) + \tfrac{\varepsilon}{2}
\le \emperr(h) + \tfrac{\varepsilon}{2}
\le \err(h) + \varepsilon$,
using representativeness, then ERM's optimality on $S$, then
representativeness again.
\end{proof}

Uniform convergence is precisely what licenses trusting the training
score of a hypothesis \emph{selected because of} its training score ---
the trap the memorizer fell into.

\begin{theorem}[Finite classes are agnostically learnable]
\label{thm:finiteagnostic}
Every finite $\Hyp$ has uniform convergence, hence is agnostically PAC
learnable by ERM, with
\begin{equation}
  m_{\Hyp}(\varepsilon,\delta)
  \;\le\;
  \Bigl\lceil \frac{2\log\bigl(2|\Hyp|/\delta\bigr)}{\varepsilon^{2}}
  \Bigr\rceil .
  \label{eq:finiteagnostic}
\end{equation}
\end{theorem}

\begin{proof}
Fix $h$. The random variables $\mathbf{1}[h(x_i)\ne y_i]$ are i.i.d.\ in
$[0,1]$ with mean $\err(h)$, so Hoeffding's inequality gives
$\Pr\bigl[|\emperr(h)-\err(h)| > \varepsilon\bigr] \le
2e^{-2m\varepsilon^2}$. A union bound over $\Hyp$ makes the failure
probability $\le 2|\Hyp| e^{-2m\varepsilon^2} \le \delta$ for $m$ as in
\eqref{eq:finiteagnostic} (at accuracy $\varepsilon/2$, whence the
constant 2). Conclude with Lemma~\ref{lem:ucerm}.
\end{proof}

Note the price of dropping realizability: $1/\varepsilon$ became
$1/\varepsilon^{2}$ --- estimating a mean to accuracy $\varepsilon$ costs
$\varepsilon^{-2}$, whereas certifying perfection costs only
$\varepsilon^{-1}$. These are the $1/m$ vs $1/\sqrt m$ rates the demo
measures.

\section{No Free Lunch: bias is necessary}

\begin{theorem}[No-Free-Lunch]
\label{thm:nfl}
Let $A$ be any learner for binary classification over a domain $\X$, and
let $m \le |\X|/2$. Then there exists a distribution $\D$, realizable by
some $f:\X\to\{0,1\}$, on which
\begin{equation}
  \Pr_{S\sim\D^m}\Bigl[\err\bigl(A(S)\bigr) \ge \tfrac18\Bigr]
  \;\ge\; \tfrac17 .
  \label{eq:nfl}
\end{equation}
\end{theorem}

(Proof idea: put uniform mass on $2m$ points; among all labelings of
those points, the sample reveals at most half, and on the unseen half any
fixed prediction is wrong for half the labelings in expectation ---
so \emph{some} labeling defeats $A$.)

\begin{corollary}
If $\X$ is infinite, the class of \emph{all} functions $\Y^{\X}$ is not
PAC learnable.
\end{corollary}

There is no universal learner: an algorithm succeeding on every problem
with a fixed budget does not exist, because for any budget an adversarial
(yet realizable!) task exists. Learning therefore always starts from
prior knowledge --- the choice of $\Hyp$. The cost structure is the
\emph{bias--complexity decomposition}: for the ERM output,
\begin{equation}
  \err(h_S)
  \;=\;
  \underbrace{\min_{h\in\Hyp}\err(h)}_{\text{approximation error}}
  \;+\;
  \underbrace{\err(h_S) - \min_{h\in\Hyp}\err(h)}_{\text{estimation error}} ,
  \label{eq:biascomplexity}
\end{equation}
richer classes lowering the first term and (by
\eqref{eq:finiteagnostic} and its VC analogue below) raising the second.
Overfitting and underfitting are the two ways of losing this trade.

\section{The VC dimension characterizes learnability}

Finiteness of $\Hyp$ was sufficient but is far from necessary ---
thresholds on $\mathbb{R}$ form an infinite class that is obviously
learnable. The right measure counts \emph{behaviors}, not hypotheses.

\begin{definition}[Shattering; VC dimension]
For $C = \{c_1,\dots,c_k\} \subseteq \X$, the \emph{restriction} of
$\Hyp$ to $C$ is
$\Hyp_C = \{(h(c_1),\dots,h(c_k)) : h\in\Hyp\} \subseteq \{0,1\}^k$.
$C$ is \emph{shattered} if $\Hyp_C = \{0,1\}^{k}$ --- every labeling of
$C$ is realized by some hypothesis. The \emph{VC dimension} $\VC(\Hyp)$
is the size of the largest shattered set (infinite if arbitrarily large
sets are shattered).
\end{definition}

\begin{example}\leavevmode
\begin{itemize}
  \item \textbf{Thresholds} $\{x\mapsto\mathbf{1}[x> a]\}$ on
    $\mathbb{R}$: any single point is shattered; for $x_1<x_2$ the
    labeling $(1,0)$ is impossible (the demo checks this exhaustively).
    $\VC = 1$.
  \item \textbf{Intervals} $\mathbf{1}[a\le x\le b]$: two points shatter;
    three fail on the labeling $(1,0,1)$. $\VC = 2$.
  \item \textbf{Axis-aligned rectangles} in $\mathbb{R}^2$: $\VC = 4$.
  \item \textbf{Halfspaces} in $\mathbb{R}^d$: $\VC = d+1$.
  \item \textbf{Finite classes}: $\VC(\Hyp) \le \log_2|\Hyp|$
    (shattering $k$ points needs $2^k$ hypotheses).
  \item \textbf{One parameter, infinite VC:}
    $\{x \mapsto \lceil\sin(\theta x)\rceil : \theta\in\mathbb{R}\}$
    shatters arbitrarily large sets --- capacity is not parameter count.
\end{itemize}
\end{example}

Why finite VC tames infinite classes: the \emph{growth function}
$\tau_{\Hyp}(m) = \max_{|C|=m} |\Hyp_C|$ counts distinguishable behaviors
on $m$ points, and once $m$ exceeds the VC dimension it stops growing
exponentially:

\begin{lemma}[Sauer--Shelah]
If $\VC(\Hyp) = d < \infty$, then
$\tau_{\Hyp}(m) \le \sum_{i=0}^{d}\binom{m}{i}$, and for $m \ge d$,
\begin{equation}
  \tau_{\Hyp}(m) \;\le\; \Bigl(\frac{em}{d}\Bigr)^{d}
  \qquad\text{--- polynomial, not } 2^m .
  \label{eq:sauer}
\end{equation}
\end{lemma}

An effective ``size'' of $(em/d)^d$ plugged into finite-class arguments
(with a symmetrization step) yields uniform convergence for any
finite-VC class --- and NFL supplies the converse: shattered sets of
unbounded size embed the unlearnable ``all functions'' problem.

\begin{theorem}[Fundamental Theorem of Statistical Learning]
\label{thm:fundamental}
For a class $\Hyp$ of binary classifiers, the following are equivalent:
\begin{enumerate}
  \item $\Hyp$ has the uniform convergence property;
  \item ERM is a successful agnostic PAC learner for $\Hyp$;
  \item $\Hyp$ is agnostically PAC learnable;
  \item $\Hyp$ is PAC learnable (realizable);
  \item $\VC(\Hyp) = d < \infty$.
\end{enumerate}
Quantitatively, with constants absorbed in $\Theta$:
\begin{equation}
  \boxed{\;
  m^{\mathrm{realizable}}_{\Hyp}(\varepsilon,\delta)
  = \Theta\!\Bigl(\frac{d + \log(1/\delta)}{\varepsilon}\Bigr),
  \qquad
  m^{\mathrm{agnostic}}_{\Hyp}(\varepsilon,\delta)
  = \Theta\!\Bigl(\frac{d + \log(1/\delta)}{\varepsilon^{2}}\Bigr).
  \;}
  \label{eq:fundamental}
\end{equation}
\end{theorem}

This is the punchline of the whole subject: \emph{learnability is a
combinatorial property of the hypothesis class} --- one number, the VC
dimension, decides yes or no and prices the data; the algorithm (ERM)
is universal among learnable classes. The demo's log-log panel shows the
rates live for thresholds ($d=1$): the realizable excess risk decays like
$1/m$, and under label noise the uniform-convergence supremum
$\sup_{h\in\Hyp}|\emperr(h)-\err(h)|$ --- representativeness over the
whole class, the quantity that costs $\varepsilon^{-2}$ --- decays like
$1/\sqrt m$.

\begin{remark}[What learnability does not promise]
(i) \emph{Computation}: the theorem counts samples, not time --- ERM can
be NP-hard even for learnable classes (e.g.\ agnostic halfspaces);
statistical and computational learnability are separate axes.
(ii) \emph{Distribution-specific help}: \eqref{eq:fundamental} is a
worst-case-over-$\D$ statement, and the demo catches this in the act:
with uniform $10\%$ label noise the measured \emph{excess risk of ERM}
decays like $1/m$, faster than the worst-case $1/\sqrt m$, even though
the uniform-convergence supremum genuinely decays like $1/\sqrt m$ on the
same data. The reason: uniform noise keeps the noise level away from
$\tfrac12$ (a \emph{margin} condition), and the linear risk penalty stops
ERM from wandering into the region where the empirical fluctuations that
drive the supremum live. Favorable distributions, margins, and
regularization all beat the worst case; that is where the rest of
learning theory (Tsybakov noise conditions, nonuniform learnability,
Rademacher complexity, stability) begins.
\end{remark}

\end{document}
